\documentclass[12pt]{article}
\usepackage[utf8]{inputenc}
\usepackage[T1]{fontenc}
\usepackage{hyperref}
\title{Saliba Language Resources for K-12 Curriculum Development}
\author{Compiled by Claude Sonnet 5.5 (claude-sonnet-5-5) for Andr\'as Kornai}
\date{}
\begin{document}
\maketitle

\section{Introduction}
Saliba (Saliba-Logea; ISO 639-3 sbe; Glottolog sali1295) is an Oceanic (Austronesian) language of Milne Bay Province, Papua New Guinea, spoken by about 2,500 people mainly on Saliba and Logea islands. Documentation exists, including a Monash University DoBeS corpus project, an online dictionary and Bible translations. Search results did not confirm any specific Saliba-medium school programme, although Milne Bay had vernacular elementary pilots from 1994 (general, not Saliba-specific).

\section{Primary Educational Resources}
No Saliba primers or textbooks were found in search results. The DoBeS project states that its materials are intended to be a source for teaching materials.
\begin{itemize}
\item Vernacular education in Papua New Guinea, UNESCO (general context, Milne Bay vernacular programmes, not Saliba-specific): \url{https://unesdoc.unesco.org/ark:/48223/pf0000146660}
\end{itemize}

\section{Linguistic and Reference Materials}
\begin{itemize}
\item Saliba online dictionary on Webonary (about 3,500 entries, compiled by R.K. Oetzel; not fully checked with speakers): \url{https://www.webonary.org/saliba/overview/introduction/}
\item Saliba-Logea documentation project (Anna Margetts, Monash University; DoBeS archive): \url{https://dobes.mpi.nl/projects/saliba/language/}
\item DoBeS poster on Saliba: \url{https://dobes.mpi.nl/posters/print/oceania-saliba-print.pdf}
\item Saliba language, Wikipedia (phonology overview): \url{https://en.wikipedia.org/wiki/Saliba_language_(Papua_New_Guinea)}
\item Saliba writing, Omniglot: \url{https://www.omniglot.com/writing/saliba.htm}
\end{itemize}

\section{Religious and Community Resources}
\begin{itemize}
\item Saliba Bible texts at eBible (PDF of the full available text): \url{https://ebible.org/pdf/sbe/sbe_all.pdf}
\item Saliba Bible at eBible (EPUB): \url{https://www.ebible.org/epub/sbe.epub}
\item IPS annual report 2017, noting the dedication of the Saliba Half New Testament on 24 May 2017: \url{https://brianc.ipsapps.org/Newsletters/png_annual-report_2017_sm.pdf}
\end{itemize}

\section{Digital and Audio Resources}
The DoBeS archive project above is described as a spoken-language corpus linking audio and video with texts, translations and annotation; access conditions were not verified.

\section{Curriculum Development Context}
Saliba has a Latin-based orthography in use in Bible translation, an online lexicon and a recorded corpus, which together are a reasonable base for developing reading materials. No published full grammar or full-scale dictionary appeared in the searches, and the claim that Saliba is used as a language of instruction was not confirmed. Next steps would be to contact the Monash/DoBeS team and Milne Bay education authorities about existing elementary materials.

\end{document}